\section{Finite-response reduction, a bounded-input estimator, and truncation}
\label{sec:transfer}
This section gives the explicit interface with nonparametric FIR estimation.
Let $\mathbf b_N(\beta)=(b_{0,\beta},\ldots,b_{N,\beta})$ and use the
maximum norm. The following deterministic assertion is independent of how
an estimated response vector was obtained.

\begin{theorem}[Response-to-physical-coefficient reduction]
\label{thm:transfer}
Suppose $\|\widehat{\mathbf b}-\mathbf b_N(\beta_0)\|_\infty\le e$ on an
event of probability at least $1-\alpha$. Let $\mathcal N_\eta$ be any
finite deterministic $\eta$-net of the compact image $\mathbf b_N(\Bcal)$,
with each net point assigned an admissible preimage. Minimize the residual
$\|\mathbf b_N(\beta)-\widehat{\mathbf b}\|_\infty$ on these finitely many
preimages, breaking ties by their fixed ordering. Then on the same event,
\begin{equation}
 d_J(\widehat\beta,\beta_0)
 \le L\{\mathcal A(2e+\eta)+\mathcal E\}.
 \label{eq:FIR-transfer}
\end{equation}
The response confidence set
$\{\beta:\|\mathbf b_N(\beta)-\widehat{\mathbf b}\|_\infty\le e\}$
is nonempty and has $d_J$ diameter at most $L(2\mathcal A e+\mathcal E)$.
Thus $\eta=e=\delta/(6L\mathcal A)$ suffices for estimation error at most
$\delta$, while $e\le\delta/(4L\mathcal A)$ suffices for confidence diameter
at most $\delta$.
\end{theorem}
\begin{proof}
Compactness and continuity give a finite net; a preimage approximating the
truth has residual at most $e+\eta$. The selected residual is no larger.
The triangle inequality bounds its response distance from the truth by
$2e+\eta$. Apply \eqref{eq:sampled-infty} and \eqref{eq:jet}. For the
confidence set use the triangle inequality between two of its elements.
The chosen $N$ ensures $L\mathcal E\le\delta/2$.
\end{proof}

\begin{theorem}[Nonreset response estimation with bounded signs]
\label{thm:response-estimator}
Fix a block count $B$, and let $M_B$ count successful blocks. For $M_B>0$ put
\[
 \widehat b_k=\frac1{aM_B}\sum_{b\le B:\,E_b}S_{b,0}Y_{b,k},\qquad
 0\le k<m,\quad
 B_*=C_*(R_0+U_*/\lambda_*),\quad v_*=\sigma^2+B_*^2.
\]
An arbitrary output may be assigned when $M_B=0$. For every $e>0$,
\begin{equation}
 \Pp\{M_B<Bp_*^m/2\text{ or }
       \|\widehat{\mathbf b}-\mathbf b_N(\beta_0)\|_\infty>e\}
 \le e^{-Bp_*^m/8}+2m\exp\{-a^2e^2Bp_*^m/(4v_*)\}.
 \label{eq:response-tail}
\end{equation}
In particular a sufficient block budget for failure at most $\alpha$ is
\begin{equation}
 B\ge\left\lceil\frac1{p_*^m}
 \max\left\{8\log\frac2\alpha,
 \frac{4v_*}{a^2e^2}\log\frac{4m}\alpha\right\}\right\rceil.
 \label{eq:response-budget}
\end{equation}
With the choice of $e,\eta$ in Theorem~\ref{thm:transfer}, this estimator
attains the same sufficient region $C_b r+p_b s<1$ as the posterior.
\end{theorem}
\begin{proof}
On a successful block, linearity gives
$Y_{b,k}=aS_{b,0}b_{k,\beta_0}+P^0_{b,k}+\xi_{b,k}$.
Condition on the starting history, success and the other block signs.
The intercept is independent of $S_{b,0}$ and has magnitude at most $B_*$:
removing the first pulse leaves a force still bounded by $U_*$. The centered
quantity $S_{b,0}Y_{b,k}-ab_{k,\beta_0}$ is conditionally sub-Gaussian with
variance proxy $v_*$. In the block filtration, therefore,
\[
 \exp\left\{\lambda\sum_{b\le B}\mathbf1_{E_b}
 (S_{b,0}Y_{b,k}-ab_{k,\beta_0})-\lambda^2v_*M_B/2\right\}
\]
is a nonnegative supermartingale. With either sign and
$\lambda=ae/v_*$, exponential Markov bounds the deviation on
$M_B\ge Bp_*^m/2$ by $2\exp\{-a^2e^2Bp_*^m/(4v_*)\}$.
The independent action patterns have success probability $p_*^m$, so the
multiplicative Chernoff count bound is $e^{-Bp_*^m/8}$. Union-bound the lags;
they are not asserted independent. This proves \eqref{eq:response-tail}
without claiming Gaussianity conditional on a terminal visit count.
For $e=\delta/(6L\mathcal A)$, $p_*^me^2=\kappa/9$. At $B\asymp n/m$
and $\alpha=n^{-u}$ for any fixed $u>0$, the sufficient budget follows from
$n\kappa/m\gg\log n$. Proposition~\ref{prop:envelope} supplies this when
$C_b r+p_b s<1$.
\end{proof}

For time-uniform bands one may instead allocate arbitrary $\eta_k>0$ summing
to one and stop each exponential supermartingale at its $l$th success. The
radii
\[
 r_{k,l}=a^{-1}\sqrt{2v_*\log(2l(l+1)/(\alpha\eta_k))/l}
\]
hold jointly for all $k$ and $l\ge1$ with probability at least $1-\alpha$.
The error allocation is $\alpha\eta_k/[l(l+1)]$. The sampling mass $w$ and
confidence allocation $\eta_k$ have different roles and are never equated.
All lags share the same success count.

\begin{proposition}[Dimension-uniform finite-section bias]
\label{prop:finite-section}
Cut the chain at site $K$ with the same diagonal, damping and interior edges,
removing only its outgoing edge. Write $b^{[K]}_{k,\beta}$ for its responses,
and $S=N\Delta+t_0$. For $K\ge1$,
\begin{equation}
 \max_{k\le N}|b_{k,\beta}-b^{[K]}_{k,\beta}|
 \le 2t_0 e^{\Lambda S}\frac{(\Lambda S)^{K+1}}{(K+1)!}
 =:\epsilon_K.\label{eq:finite-bias}
\end{equation}
The same stability constants work for every cut chain, and
\begin{equation}
 \sum_{k\ge0}|b^{[K]}_{k,\beta}|
 \le\frac{C_*t_0e^{\Lambda t_0}}{1-e^{-\lambda_*\Delta}}
 \label{eq:finite-uniform}
\end{equation}
is independent of $K$. To make $\epsilon_K\le\epsilon$, it suffices that
\[
 K+1\ge\max\left\{2,2e\Lambda S,
  \frac{\log(2t_0/\epsilon)+\Lambda S}{\log2}\right\}.
\]
Hence at $N=O(\log n)$ and a polynomially small response error budget,
$K=O(\log n)$ suffices uniformly in the infinite tail.
\end{proposition}
\begin{proof}
Boundary generator words agree through order $K$; a path leaving the cut and
returning cannot be shorter. The difference of exponential series at time
$k\Delta+u$ is bounded by
$2e^{\Lambda S}(\Lambda S)^{K+1}/(K+1)!$. Integrate over $0\le u\le t_0$.
The quadratic-form bounds on the Jacobi matrix and the energy proof of
Lemma~\ref{lem:stability} are unchanged by the cut. Applying stability to
$U_K^kH_KB$ and summing the geometric series proves
\eqref{eq:finite-uniform}. Finally $q!\ge(q/e)^q$ and $q\ge2e\Lambda S$
give $(\Lambda S)^q/q!\le2^{-q}$.
\end{proof}

\subsection{What a finite-response theorem supplies, and what remains to prove}
The uniform envelope \eqref{eq:finite-uniform} shows why latent dimension need
not force a loss in a response-level system norm. Equation~\eqref{eq:finite-bias}
quantifies the separate truncation bias instead of setting it to zero.
For example, at fixed $r,s$, the tolerance
$e=\delta/(6L\mathcal A)$ has $\log(1/e)=O(\log n)$, so a logarithmic-size
finite section resolves the required response accuracy. If an FIR theorem
under a compatible experiment supplies $\ell^2$ error $e$, it also supplies
the maximum-norm error used in Theorem~\ref{thm:transfer}; no realization or
similarity-coordinate recovery is required at this stage.

Compatibility is substantive. The Gaussian inputs of Algorithm 1 in
\cite{SRD2020} are not uniformly bounded, and its printed initial-state and
process-noise hypotheses are not the present experiment. Simply citing that
algorithm would not prove our bounded-input, retained-nuisance claim.
Theorem~\ref{thm:response-estimator} supplies the needed front end directly.
Conversely, there is no claim that this front end replaces the broader
nonparametric results of \cite{Gold1998,SRD2020}. The transfer theorem exposes
exactly how such an estimate can be reused when its hypotheses are verified.
The sampled-channel inverse, the tail-uniform physical conditioning, and the
complete-posterior transfer are the additional mathematical steps. A cut-chain
calculation is only a computational approximation: its removed edge is not
silently treated as an admissible positive infinite tail.
